%% 2-19 プログラムで計算しよう（第11回：ページ制限なし。変数の中身を表で追う・なぜ 0.1 がずれるのかを1段ずつ）
%% この章の部品（頭文字 \bkiit…）
%%   プログラムの枠  \begin{bkiitcode}[題] … \end{bkiitcode}（中の改行と字下げの空白は、そのまま出る）
%%   実行結果の枠    \begin{bkiitout} … \end{bkiitout}
%%   ブロックの絵    \bkiitblock{(x,y)}{はば}{色}{文字}、くり返しの枠 \bkiitloop{(x,y)}{はば}{中の高さ}{色}{文字}
\makeatletter
{\catcode`\ =\active\gdef\bkiit@obey{\catcode`\ =\active\def {\leavevmode\nobreak\ }\obeylines\parindent=0pt\parskip=0pt}}
\newtcolorbox{bkiitcode}[1][Python]{enhanced, colback=bkPaper, colframe=bkNavy!60, boxrule=0.7pt, arc=1.5mm,
  left=3mm, right=3mm, top=1.2mm, bottom=1.2mm, before skip=3pt, after skip=3pt,
  fontupper=\ttfamily\gtfamily\small, before upper={\bkiit@obey\baselineskip=14.5pt},
  attach boxed title to top right={xshift=-3mm, yshift*=-\tcboxedtitleheight/2},
  boxed title style={colback=bkNavy!60, colframe=bkNavy!60, arc=1mm, boxrule=0pt, left=1.5mm, right=1.5mm, top=0.3mm, bottom=0.3mm},
  fonttitle=\sffamily\bfseries\scriptsize, coltitle=white, title={#1}}
\newtcolorbox{bkiitout}[1][\R{実行}{じっ|こう}\R{結果}{けっ|か}]{enhanced, colback=bkInk!85, colframe=bkInk!85, boxrule=0.7pt, arc=1.5mm,
  left=3mm, right=3mm, top=1.2mm, bottom=1.2mm, before skip=3pt, after skip=3pt, coltext=white,
  fontupper=\ttfamily\gtfamily\small, before upper={\bkiit@obey\baselineskip=14.5pt},
  attach boxed title to top right={xshift=-3mm, yshift*=-\tcboxedtitleheight/2},
  boxed title style={colback=bkGray, colframe=bkGray, arc=1mm, boxrule=0pt, left=1.5mm, right=1.5mm, top=0.3mm, bottom=0.3mm},
  fonttitle=\sffamily\bfseries\scriptsize, coltitle=white, title={#1}}
% ブロック（上にくぼみ、下に出っぱり）。(x,y) は左上
\newcommand\bkiitblock[4]{\begin{scope}[shift={#1}]
  \filldraw[fill=#3, draw=#3!70!black, line width=0.5pt, rounded corners=0.6mm]
    (0,0) -- (3,0) -- (4,-1) -- (7,-1) -- (8,0) -- (#2,0) -- (#2,-7) -- (8,-7) -- (7,-8) -- (4,-8) -- (3,-7) -- (0,-7) -- cycle;
  \node[anchor=west, font=\sffamily\small, text=white, inner sep=0pt] at (2.5,-3.7) {#4};
  \end{scope}}
% くり返し・もしの枠（Cの形）。中の高さ #3、上の帯に文字
\newcommand\bkiitloop[5]{\begin{scope}[shift={#1}]
  \filldraw[fill=#4, draw=#4!70!black, line width=0.5pt, rounded corners=0.6mm]
    (0,0) -- (3,0) -- (4,-1) -- (7,-1) -- (8,0) -- (#2,0) -- (#2,-7)
    -- (12,-7) -- (11,-8) -- (8,-8) -- (7,-7) -- (4,-7)
    -- (4,-7-#3) -- (7,-7-#3) -- (8,-8-#3) -- (11,-8-#3) -- (12,-7-#3) -- (#2,-7-#3)
    -- (#2,-11-#3) -- (8,-11-#3) -- (7,-12-#3) -- (4,-12-#3) -- (3,-11-#3) -- (0,-11-#3) -- cycle;
  \node[anchor=west, font=\sffamily\small, text=white, inner sep=0pt] at (2.5,-3.7) {#5};
  \end{scope}}
% もし～なら／でなければ（口が2つ）  \bkiitifelse{(x,y)}{はば}{上の口の高さ}{下の口の高さ}{色}{上の帯の文字}{中の帯の文字}
\newcommand\bkiitifelse[7]{\begin{scope}[shift={#1}]
  \filldraw[fill=#5, draw=#5!70!black, line width=0.5pt, rounded corners=0.6mm]
    (0,0) -- (3,0) -- (4,-1) -- (7,-1) -- (8,0) -- (#2,0) -- (#2,-7)
    -- (12,-7) -- (11,-8) -- (8,-8) -- (7,-7) -- (4,-7)
    -- (4,-7-#3) -- (7,-7-#3) -- (8,-8-#3) -- (11,-8-#3) -- (12,-7-#3) -- (#2,-7-#3)
    -- (#2,-14-#3) -- (12,-14-#3) -- (11,-15-#3) -- (8,-15-#3) -- (7,-14-#3) -- (4,-14-#3)
    -- (4,-14-#3-#4) -- (7,-14-#3-#4) -- (8,-15-#3-#4) -- (11,-15-#3-#4) -- (12,-14-#3-#4) -- (#2,-14-#3-#4)
    -- (#2,-18-#3-#4) -- (8,-18-#3-#4) -- (7,-19-#3-#4) -- (4,-19-#3-#4) -- (3,-18-#3-#4) -- (0,-18-#3-#4) -- cycle;
  \node[anchor=west, font=\sffamily\small, text=white, inner sep=0pt] at (2.5,-3.7) {#6};
  \node[anchor=west, font=\sffamily\small, text=white, inner sep=0pt] at (2.5,-10.7-#3) {#7};
  \end{scope}}
% 変数の箱  \bkiitbox{(x,y)}{名前}{中身}{色}  (x,y) は箱のまん中
\newcommand\bkiitbox[4]{\begin{scope}[shift={#1}]
  \filldraw[fill=white, draw=#4, line width=1pt, rounded corners=0.8mm] (-7,-4.5) rectangle (7,4.5);
  \node[font=\sffamily\bfseries\large, text=bkInk] at (0,0) {#3};
  \node[anchor=south, font=\ttfamily\small, text=#4!80!black, inner sep=0pt] at (0,6) {#2};
  \end{scope}}
% 変数の中身を追う表（行｜変数…）
\NewTblrEnviron{bkiittrace}
\SetTblrInner[bkiittrace]{rowsep=0.6pt, colsep=3pt, hlines={0.4pt,bkLine}, vlines={0.4pt,bkLine},
  row{1}={bg=bkNavy!12, font=\sffamily\bfseries\footnotesize}, cells={font=\small}}
\makeatother

\bkchapter{2}{19}{プログラムで\R{計算}{けい|さん}しよう}{\bklevel{2}}{プログラムで\R{計算}{けい|さん}しよう}
\bkmichishirube{2-9,2-17,2-18}{5-2,6-2}

\begin{goal}
\begin{itemize}[label={\color{bkTheme}●}, leftmargin=1.6\zw, itemsep=0pt, topsep=1pt]
\item \bklevel{2}\enspace \R{変数}{へん|すう}と\R{代入}{だい|にゅう}を\R{使}{つか}った\R{短}{みじか}いプログラムを\R{読}{よ}んで、\R{結果}{けっ|か}を\R{予想}{よ|そう}できる。
\item \bklevel{2}\enspace くり\R{返}{かえ}しと「もし～なら」を\R{使}{つか}ったプログラムを\R{読}{よ}める。
\item \bklevel{2}\enspace コンピュータの\R{小数}{しょう|すう}に\R{小}{ちい}さなずれが\R{出}{で}るわけと、その\R{対策}{たい|さく}を\R{説明}{せつ|めい}できる。
\end{itemize}
\end{goal}

\begin{hatena}
\noindent\begin{minipage}[c]{0.56\linewidth}
コンピュータに「0.1を10\R{回}{かい}たして」とたのみました。\R{答}{こた}えは1…のはずです。ところが、\R{画面}{が|めん}に\R{出}{で}たのは\R{右}{みぎ}の\R{数}{かず}でした。

コンピュータが、\R{計算}{けい|さん}をまちがえたのかな？
\end{minipage}\hfill
\begin{minipage}[c]{0.41\linewidth}\centering
\begin{tikzpicture}[x=1mm, y=1mm]
  \fill[bkInk!85, rounded corners=2mm] (0,0) rectangle (48,24);
  \fill[bkGray!60, rounded corners=2mm] (0,19) rectangle (48,24);
  \fill[bkInk!85] (0,19) rectangle (48,20.5);
  \foreach \c [count=\k] in {bkRed,bkYellow,bkGreen} \fill[\c] (\k*4,21.7) circle (1);
  \node[anchor=west, font=\ttfamily\small, text=white] at (3,13.5) {0.1を10\R{回}{かい}たすと…};
  \node[anchor=west, font=\ttfamily\small, text=bkYellow] at (3,6) {0.9999999999999999};
\end{tikzpicture}
\end{minipage}
\end{hatena}

\noindent 0.1を10\R{回}{かい}たすのは、\R{電卓}{でん|たく}や\R{手計算}{て|けい|さん}なら\R{誰}{だれ}でもできる\R{計算}{けい|さん}です。ふしぎなずれのわけは、コンピュータの\R{数}{かず}の\R{覚}{おぼ}え\R{方}{かた}（2-17）にあります。この\R{章}{しょう}では、コンピュータに\R{計算}{けい|さん}の\R{手順}{て|じゅん}を\R{伝}{つた}える\textbf{プログラム}の\R{読}{よ}み\R{方}{かた}を\R{学}{まな}び、このなぞを\R{解}{と}きます。\bkidx{ぷろぐらむ}{プログラム}

%% ---------------- やってみよう・思い出そう ----------------
\begin{yattemiyou}[ロボットになって、\R{命令}{めい|れい}どおりに\R{動}{うご}こう]
ふたりで\R{組}{く}みます。ひとりが\R{命令}{めい|れい}を\R{書}{か}き、もうひとりがロボットになって、\R{書}{か}かれたとおりに\R{動}{うご}きます。\R{使}{つか}える\R{命令}{めい|れい}は「\R{前}{まえ}へ\R{何歩}{なん|ぽ}」と「\R{右}{みぎ}を\R{向}{む}く」の2つだけです。

\noindent\begin{minipage}[c]{0.56\linewidth}
\begin{enumerate}[label={\bkcircled[bkGreen]{\arabic*}}, leftmargin=2.4\zw, itemsep=1pt, topsep=2pt]
\item 1\R{辺}{ぺん}2\R{歩}{ほ}の\R{正方形}{せい|ほう|けい}をまわって\R{元}{もと}の\R{場所}{ば|しょ}・\R{元}{もと}の\R{向}{む}きにもどる\R{命令}{めい|れい}を、ノートに\R{順番}{じゅん|ばん}に\R{書}{か}こう。
\item ロボットに、\R{書}{か}いた\R{順番}{じゅん|ばん}どおりに\R{動}{うご}いてもらおう。\R{命令}{めい|れい}にないことはしない。
\item \R{同}{おな}じ\R{命令}{めい|れい}の\R{組}{く}が、\R{何回}{なん|かい}\R{出}{で}てきたかな？\quad もっと\R{短}{みじか}く\R{書}{か}けないかな？
\end{enumerate}
\end{minipage}\hfill
\begin{minipage}[c]{0.4\linewidth}\centering
\begin{tikzpicture}[x=1mm, y=1mm]
  \foreach \i in {0,...,2} {\draw[bkLine, line width=0.5pt] (\i*10,0) -- (\i*10,20); \draw[bkLine, line width=0.5pt] (0,\i*10) -- (20,\i*10);}
  \draw[bkGreen, line width=1.4pt, -{Stealth[length=5pt]}] (0,0) -- (0,19.2);
  \draw[bkGreen, line width=1.4pt, -{Stealth[length=5pt]}] (0,20) -- (19.2,20);
  \draw[bkGreen, line width=1.4pt, -{Stealth[length=5pt]}] (20,20) -- (20,0.8);
  \draw[bkGreen, line width=1.4pt, -{Stealth[length=5pt]}] (20,0) -- (0.8,0);
  \fill[bkGreen] (0,0) circle (1.2);
  \bkkid{(-9,-1)}{12}{bkBlue}{bkHair}
  \node[anchor=north, font=\sffamily\footnotesize, text=bkInk] at (6,-3.5) {スタート（\R{上}{うえ}\R{向}{む}き）};
\end{tikzpicture}
\end{minipage}
\end{yattemiyou}

\noindent ①の\R{命令}{めい|れい}は、「\R{前}{まえ}へ2\R{歩}{ほ}、\R{右}{みぎ}を\R{向}{む}く、\R{前}{まえ}へ2\R{歩}{ほ}、\R{右}{みぎ}を\R{向}{む}く、…」と8\R{行}{ぎょう}になります。「\R{前}{まえ}へ2\R{歩}{ほ}・\R{右}{みぎ}を\R{向}{む}く」の\R{組}{く}が4\R{回}{かい}\R{出}{で}てくるので、「この2つを4\R{回}{かい}くり\R{返}{かえ}す」と\R{書}{か}けば、3\R{行}{ぎょう}ですみます。ロボットは、\R{書}{か}かれていないことは\R{何}{なに}もしません。\R{命令}{めい|れい}の\R{順番}{じゅん|ばん}をまちがえると、ちがう\R{場所}{ば|しょ}に\R{行}{い}ってしまいます。コンピュータへの\R{命令}{めい|れい}も、これと\R{同}{おな}じです。

\begin{omoidasou}
\bkrecall{2-9}{$\dfrac{1}{3}=0.333\cdots$ のように、\R{分数}{ぶん|すう}を\R{小数}{しょう|すう}にすると、\R{終}{お}わらずにくり\R{返}{かえ}すことがある。}
\bkrecall{2-17}{2\R{進数}{しん|すう}の\R{位}{くらい}は、\R{右}{みぎ}から1・2・4・8…。\R{大}{おお}きい\R{位}{くらい}から、\R{使}{つか}えたらひいて2\R{進数}{しん|すう}にする。}
\bkrecall{2-7}{$2^{-1}=\dfrac{1}{2}$、$2^{-2}=\dfrac{1}{4}$、$2^{-3}=\dfrac{1}{8}$。$2^{10}=1024$ はほぼ1000。}
\bkrecall{2-18}{\R{表計算}{ひょう|けい|さん}では、\R{上}{うえ}の\R{行}{ぎょう}の\R{結果}{けっ|か}を\R{使}{つか}って\R{下}{した}の\R{行}{ぎょう}を\R{計算}{けい|さん}し、\R{式}{しき}をコピーしてくり\R{返}{かえ}せる。×は *。}
\bkrecall{2-16}{$x=x+1$ を\R{移項}{い|こう}すると $0=1$ となり、\R{方程式}{ほう|てい|しき}としては\R{解}{かい}がない。}
\end{omoidasou}

%% ---------------- 数学の道具 ----------------
\bksection{doguu}{\R{数学}{すう|がく}の\R{道具}{どう|ぐ}}

\bksubsection{1}{プログラムと\R{変数}{へん|すう}・\R{代入}{だい|にゅう}}

\leavevmode\label{bk:2-19:dainyu}コンピュータへの\R{命令}{めい|れい}を\R{順番}{じゅん|ばん}にならべたものを\textbf{プログラム}といいます。ブロックをつないで\R{書}{か}く\R{方法}{ほう|ほう}（Scratchなど）と、\R{文字}{も|じ}で\R{書}{か}く\R{方法}{ほう|ほう}（Pythonなど）があります。\R{同}{おな}じプログラムを、2つの\R{書}{か}き\R{方}{かた}でならべてみましょう。

\noindent\begin{minipage}[c]{0.45\linewidth}\centering
\begin{tikzpicture}[x=1mm, y=1mm]
  \bkiitblock{(0,0)}{46}{bkOrange}{x を 3 にする}
  \bkiitblock{(0,-8)}{46}{bkOrange}{y を x×2＋1 にする}
  \bkiitblock{(0,-16)}{46}{bkPurple}{y を\R{表示}{ひょう|じ}する}
\end{tikzpicture}
\end{minipage}\hfill
\begin{minipage}[c]{0.52\linewidth}
\begin{bkiitcode}
x = 3
y = x * 2 + 1
print(y)
\end{bkiitcode}
\begin{bkiitout}
7
\end{bkiitout}
\end{minipage}

\par\addvspace{4pt}
プログラムは、\R{上}{うえ}から\R{順}{じゅん}に1\R{行}{ぎょう}ずつ\R{実行}{じっ|こう}されます。x や y は、\R{数}{かず}を\R{入}{い}れておく\R{名前}{な|まえ}つきの\R{箱}{はこ}で、\textbf{\R{変数}{へん|すう}}といいます（2-11の「どんな\R{数}{かず}でも\R{入}{はい}る\R{箱}{はこ}」と\R{同}{おな}じ）。\bkidx{へんすう}{\R{変数}{へん|すう}}\mbox{\texttt{=} は}「\R{右}{みぎ}を\R{計算}{けい|さん}して、\R{答}{こた}えを\R{左}{ひだり}の\R{箱}{はこ}に\R{入}{い}れる」という\R{命令}{めい|れい}で、\textbf{\R{代入}{だい|にゅう}}といいます。\texttt{*} はかけ\R{算}{ざん}（2-18と\R{同}{おな}じ）、\texttt{print} は「\R{表示}{ひょう|じ}する」という\R{命令}{めい|れい}です。

1\R{行}{ぎょう}ずつ、\R{箱}{はこ}の\R{中身}{なか|み}がどうなるかを\R{見}{み}てみましょう。

\begin{bkfigure}
\begin{tikzpicture}[x=1mm, y=1mm]
  \node[anchor=west, font=\sffamily\small, text=bkInk] at (0,0) {① \texttt{x = 3}};
  \bkiitbox{(44,0)}{x}{3}{bkOrange}
  \bkiitbox{(64,0)}{y}{}{bkBlue}
  \node[anchor=west, font=\sffamily\footnotesize, text=bkInk] at (76,0) {x の\R{箱}{はこ}に3を\R{入}{い}れる};
  \node[anchor=west, font=\sffamily\small, text=bkInk] at (0,-18) {② \texttt{y = x * 2 + 1}};
  \bkiitbox{(44,-18)}{x}{3}{bkOrange}
  \bkiitbox{(64,-18)}{y}{7}{bkBlue}
  \node[anchor=west, font=\sffamily\footnotesize, text=bkInk, align=left] at (76,-18) {x の\R{中身}{なか|み}3を\R{使}{つか}って\\$3\times2+1=7$ を y へ};
  \node[anchor=west, font=\sffamily\small, text=bkInk] at (0,-36) {③ \texttt{print(y)}};
  \bkiitbox{(44,-36)}{x}{3}{bkOrange}
  \bkiitbox{(64,-36)}{y}{7}{bkBlue}
  \node[anchor=west, font=\sffamily\footnotesize, text=bkInk] at (76,-36) {y の\R{中身}{なか|み}7を\R{画面}{が|めん}に\R{出}{だ}す};
\end{tikzpicture}
\bkfigcaption{1\R{行}{ぎょう}\R{実行}{じっ|こう}するごとの\R{箱}{はこ}の\R{中身}{なか|み}。②では、かけ\R{算}{ざん}を\R{先}{さき}に\R{計算}{けい|さん}する（2-11）。}
\end{bkfigure}

この\R{箱}{はこ}の\R{中身}{なか|み}の\R{変}{か}わり\R{方}{かた}は、\R{表}{ひょう}にするとわかりやすくなります。1\R{行}{ぎょう}\R{実行}{じっ|こう}するごとに、\R{表}{ひょう}の1\R{行}{ぎょう}をうめていきます。\R{空}{から}の\R{箱}{はこ}は「－」と\R{書}{か}きます。

\begin{bkcenter}
\begin{bkiittrace}{colspec={Q[l,m,wd=36mm] Q[c,m,wd=12mm] Q[c,m,wd=12mm] Q[l,m,wd=30mm]}}
\R{実行}{じっ|こう}した\R{行}{ぎょう} & x & y & \R{表示}{ひょう|じ} \\
\texttt{x = 3} & 3 & － & \\
\texttt{y = x * 2 + 1} & 3 & 7 & \\
\texttt{print(y)} & 3 & 7 & 7 \\
\end{bkiittrace}
\bktabcaption{\R{変数}{へん|すう}の\R{中身}{なか|み}を\R{追}{お}う\R{表}{ひょう}。プログラムを\R{読}{よ}むときの\R{強}{つよ}い\R{味方}{み|かた}。}
\end{bkcenter}

\bksubsection{2}{\texttt{x = x + 1} は\R{方程式}{ほう|てい|しき}ではない}

\leavevmode\label{bk:2-19:xx}では、\texttt{x = x + 1} は\R{何}{なに}をするのでしょう。\R{方程式}{ほう|てい|しき} $x=x+1$ なら、$x$ を\R{移項}{い|こう}すると $0=1$ になり、\R{解}{かい}はありません（2-16）。でも、プログラムの \texttt{=} は「\R{等}{ひと}しい」ではなく、\R{代入}{だい|にゅう}です。x に3が\R{入}{はい}っているとき、コンピュータは\R{次}{つぎ}の\R{順}{じゅん}で\R{動}{うご}きます。

\begin{enumerate}[label={\bkcircled{\arabic*}}, leftmargin=2.4\zw, itemsep=0pt, topsep=2pt]
\item \R{右}{みぎ}がわの \texttt{x + 1} を\R{見}{み}る。x の\R{箱}{はこ}の\R{中身}{なか|み}は3。
\item $3+1=4$ を\R{計算}{けい|さん}する。
\item \R{答}{こた}えの4を、\R{左}{ひだり}がわの x の\R{箱}{はこ}に\R{入}{い}れる。\R{前}{まえ}に\R{入}{はい}っていた3は\R{消}{き}えて、4に\R{入}{い}れかわる。
\end{enumerate}

\begin{bkfigure}
\begin{tikzpicture}[x=1mm, y=1mm]
  \node[draw=bkOrange, line width=1pt, fill=white, minimum width=14mm, minimum height=9mm, font=\sffamily\bfseries\large] (b1) at (0,0) {3};
  \node[anchor=south, font=\sffamily\footnotesize, text=bkInk, inner sep=0pt] at ([yshift=1.5mm]b1.north) {\R{前}{まえ}の \texttt{x}};
  \node[draw=bkOrange, line width=1pt, fill=white, minimum width=14mm, minimum height=9mm, font=\sffamily\bfseries\large] (b2) at (80,0) {4};
  \node[anchor=south, font=\sffamily\footnotesize, text=bkInk, inner sep=0pt] at ([yshift=1.5mm]b2.north) {あとの \texttt{x}};
  \node[bk node, font=\sffamily\small] (c) at (40,0) {\R{右}{みぎ}の \texttt{x + 1} を\R{計算}{けい|さん}：$3+1=4$};
  \draw[bk link] (b1) -- (c);
  \draw[bk link] (c) -- (b2);
\end{tikzpicture}
\bkfigcaption{\texttt{x = x + 1}：\R{今}{いま}の x に1をたして、x に\R{入}{い}れなおす。}
\end{bkfigure}

\R{黒板}{こく|ばん}に\R{書}{か}いた\R{点数}{てん|すう}を\R{消}{け}して、1\R{点}{てん}ふやした\R{数}{かず}を\R{書}{か}きなおすのと\R{同}{おな}じです。だから \texttt{x = x + 1} は「x を1ふやす」という\R{命令}{めい|れい}になります。2-18の \texttt{=B2+500} が「\R{上}{うえ}のセルに500をたしてここに\R{入}{い}れる」だったのと、よく\R{似}{に}ています。プログラムで「\R{等}{ひと}しいか」を\R{比}{くら}べたいときは、\texttt{=} ではなく \texttt{==} を\R{使}{つか}います。

\bksubsection{3}{くり\R{返}{かえ}し}

\leavevmode\label{bk:2-19:kurikaesi}2-18の\R{貯金}{ちょ|きん}（3000\R{円}{えん}から\R{毎月}{まい|つき}500\R{円}{えん}）を、12か\R{月分}{げつ|ぶん}\R{計算}{けい|さん}するプログラムです。

\noindent\begin{minipage}[c]{0.45\linewidth}\centering
\begin{tikzpicture}[x=1mm, y=1mm]
  \bkiitblock{(0,0)}{52}{bkOrange}{money を 3000 にする}
  \bkiitloop{(0,-8)}{52}{8}{bkGreen}{12\R{回}{かい}くり\R{返}{かえ}す}
  \bkiitblock{(4,-15)}{48}{bkOrange}{money を money＋500 に}
  \bkiitblock{(0,-27)}{52}{bkPurple}{money を\R{表示}{ひょう|じ}する}
\end{tikzpicture}
\end{minipage}\hfill
\begin{minipage}[c]{0.52\linewidth}
\begin{bkiitcode}
money = 3000
for i in range(12):
    money = money + 500
print(money)
\end{bkiitcode}
\begin{bkiitout}
9000
\end{bkiitout}
\end{minipage}

\par\addvspace{4pt}
\begin{itemize}[leftmargin=1.2\zw, itemsep=1pt, topsep=2pt, label={\color{bkTheme}●}]
\item \texttt{for i in range(12):} は「\R{下}{した}の\R{命令}{めい|れい}を12\R{回}{かい}くり\R{返}{かえ}す」という\R{意味}{い|み}です。i は\R{回数}{かい|すう}を\R{数}{かぞ}える\R{係}{かかり}の\R{変数}{へん|すう}で、0、1、2、…、11と12\R{回}{かい}\R{変}{か}わります（0から\R{数}{かぞ}えるので、\R{最後}{さい|ご}は11）。
\item くり\R{返}{かえ}すのは、\R{右}{みぎ}へ\R{下}{さ}げて（4\R{文字}{も|じ}\R{分}{ぶん}あけて）\R{書}{か}いた\R{行}{ぎょう}だけです。Scratchの\R{絵}{え}で、くり\R{返}{かえ}しの\R{枠}{わく}の\R{中}{なか}に\R{入}{はい}っているブロックにあたります。
\item \texttt{print(money)} は\R{下}{さ}げていないので、くり\R{返}{かえ}しが\R{終}{お}わってから1\R{回}{かい}だけ\R{実行}{じっ|こう}されます。
\end{itemize}

\R{変数}{へん|すう}の\R{表}{ひょう}で\R{追}{お}ってみましょう。

\begin{bkcenter}
\begin{bkiittrace}{colspec={Q[l,m,wd=50mm] Q[c,m,wd=12mm] Q[c,m,wd=16mm] Q[l,m,wd=16mm]}}
\R{実行}{じっ|こう}した\R{行}{ぎょう} & i & money & \R{表示}{ひょう|じ} \\
\texttt{money = 3000} & － & 3000 & \\
1\R{回目}{かい|め}：\texttt{money = money + 500} & 0 & 3500 & \\
2\R{回目}{かい|め}：\texttt{money = money + 500} & 1 & 4000 & \\
3\R{回目}{かい|め}：\texttt{money = money + 500} & 2 & 4500 & \\
$\vdots$ & $\vdots$ & $\vdots$ & \\
12\R{回目}{かい|め}：\texttt{money = money + 500} & 11 & 9000 & \\
\texttt{print(money)} & 11 & 9000 & 9000 \\
\end{bkiittrace}
\bktabcaption{くり\R{返}{かえ}すたびに、money が500ずつふえる。2-18の\R{表計算}{ひょう|けい|さん}のB\R{列}{れつ}と\R{同}{おな}じ\R{数}{かず}がならぶ。}
\end{bkcenter}

\R{表計算}{ひょう|けい|さん}では、\R{式}{しき}をコピーした\R{行}{ぎょう}の\R{数}{かず}だけ\R{計算}{けい|さん}がくり\R{返}{かえ}されました。プログラムでは、\R{同}{おな}じ1つの\R{箱}{はこ}の\R{中身}{なか|み}を、\texttt{range} の\R{回数}{かい|すう}だけ\R{書}{か}きかえていきます。\R{答}{こた}えは、\R{式}{しき}で\R{計算}{けい|さん}した $3000+500\times12=3000+6000=9000$ と\R{同}{おな}じです。

\bksubsection{4}{もし～なら：\R{条件分岐}{じょう|けん|ぶん|き}}

\leavevmode\label{bk:2-19:moshi}12か\R{月後}{げつ|ご}に10000\R{円}{えん}たまったかどうかで、\R{表示}{ひょう|じ}することばを\R{変}{か}えます。

\noindent\begin{minipage}[c]{0.45\linewidth}\centering
\begin{tikzpicture}[x=1mm, y=1mm]
  \bkiitifelse{(0,0)}{52}{8}{8}{bkGreen}{もし money≧10000 なら}{でなければ}
  \bkiitblock{(4,-7)}{40}{bkPurple}{「たまった」と\R{表示}{ひょう|じ}}
  \bkiitblock{(4,-22)}{40}{bkPurple}{「まだ」と\R{表示}{ひょう|じ}}
\end{tikzpicture}
\par\vspace{2pt}{\raggedright\footnotesize\sffamily ※ Scratch には≧のブロックがないので、「money＞9999」のように\R{書}{か}きます。\par}
\end{minipage}\hfill
\begin{minipage}[c]{0.52\linewidth}
\begin{bkiitcode}
if money >= 10000:
    print("たまった")
else:
    print("まだ")
\end{bkiitcode}
\begin{bkiitout}
まだ
\end{bkiitout}
\end{minipage}

\par\addvspace{4pt}
\texttt{if} の\R{後}{うし}ろの\R{条件}{じょう|けん}（\texttt{money >= 10000}）が\R{正}{ただ}しければ、すぐ\R{下}{した}の\R{下}{さ}げた\R{行}{ぎょう}を\R{実行}{じっ|こう}します。\R{正}{ただ}しくなければ、\texttt{else} の\R{下}{した}の\R{行}{ぎょう}を\R{実行}{じっ|こう}します。このように、\R{条件}{じょう|けん}が\R{正}{ただ}しいかどうかで\R{実行}{じっ|こう}する\R{命令}{めい|れい}を\R{変}{か}えることを\textbf{\R{条件分岐}{じょう|けん|ぶん|き}}といいます。\bkidx{じょうけんぶんき}{\R{条件分岐}{じょう|けん|ぶん|き}}\R{道}{みち}が2つに\R{分}{わ}かれていて、\R{条件}{じょう|けん}で\R{進}{すす}む\R{道}{みち}を\R{選}{えら}ぶイメージです。

money は9000です。$9000\geqq10000$ は\R{正}{ただ}しくない（9000は10000より\R{小}{ちい}さい）ので、\texttt{else} の\R{下}{した}の「まだ」が\R{表示}{ひょう|じ}されます。

\begin{bkcenter}
\begin{bktable}{colspec={Q[c,m,wd=14mm] Q[c,m,wd=14mm] Q[l,m,wd=58mm]}, rowsep=0.8pt, colsep=4pt, cells={font=\small}, row{1}={font=\sffamily\bfseries\small}, column{1}={font=\ttfamily\small}}
{\sffamily プログラム} & \R{算数}{さん|すう} & \R{意味}{い|み} \\
>= & $\geqq$ & \R{左}{ひだり}が\R{右}{みぎ}\R{以上}{い|じょう}（2-2） \\
<= & $\leqq$ & \R{左}{ひだり}が\R{右}{みぎ}\R{以下}{い|か} \\
> & $>$ & \R{左}{ひだり}が\R{右}{みぎ}より\R{大}{おお}きい \\
< & $<$ & \R{左}{ひだり}が\R{右}{みぎ}より\R{小}{ちい}さい（\R{未満}{み|まん}） \\
== & $=$ & \R{左}{ひだり}と\R{右}{みぎ}が\R{等}{ひと}しい \\
\end{bktable}
\bktabcaption{くらべる\R{記号}{き|ごう}。キーボードに≧・≦がないので、2\R{文字}{も|じ}で\R{書}{か}く。}
\end{bkcenter}

\bksubsection{5}{0.1を10\R{回}{かい}たしても、1にならない}

\leavevmode\label{bk:2-19:zure}\bkhatenaref を\R{確}{たし}かめましょう。1\R{回}{かい}たすごとに x を\R{表示}{ひょう|じ}します。\texttt{print(x)} を\R{右}{みぎ}へ\R{下}{さ}げてあるので、くり\R{返}{かえ}しの\R{中}{なか}で\R{毎回}{まい|かい}\R{表示}{ひょう|じ}されます。

\noindent\begin{minipage}[c]{0.36\linewidth}
\begin{bkiitcode}
x = 0
for i in range(10):
    x = x + 0.1
    print(x)
\end{bkiitcode}
\end{minipage}\hfill
\begin{minipage}[c]{0.62\linewidth}\centering
\begin{tikzpicture}[x=1mm, y=1mm, outl/.style={anchor=north west, align=left, text=white, font=\ttfamily\footnotesize, inner sep=0pt, execute at begin node={\baselineskip=13pt}}]
  \fill[bkInk!85, rounded corners=1.5mm] (0,0) rectangle (84,-30);
  \node[fill=bkGray, rounded corners=1mm, text=white, font=\sffamily\bfseries\scriptsize, inner xsep=1.5mm, inner ysep=0.6mm] at (72,0) {\R{実行}{じっ|こう}\R{結果}{けっ|か}};
  \node[outl] at (3,-4.5) {0.1\\0.2\\{\color{bkYellow}0.30000000000000004}\\0.4\\0.5};
  \node[outl] at (45,-4.5) {0.6\\0.7\\{\color{bkYellow}0.7999999999999999}\\{\color{bkYellow}0.8999999999999999}\\{\color{bkYellow}0.9999999999999999}};
\end{tikzpicture}
\par{\footnotesize\sffamily \R{左}{ひだり}の\R{列}{れつ}のあとに\R{右}{みぎ}の\R{列}{れつ}が\R{出}{で}る。\R{黄色}{き|いろ}は、ずれた\R{数}{かず}。\par}
\end{minipage}

\par\addvspace{4pt}
3\R{回目}{かい|め}から\R{小}{ちい}さなずれが\R{見}{み}えます。わけは、コンピュータが\R{数}{かず}を\textbf{2\R{進数}{しん|すう}で、\R{決}{き}まったけた\R{数}{すう}だけ}\R{覚}{おぼ}えているからです（2-17）。

\textbf{2\R{進数}{しん|すう}の\R{小数}{しょう|すう}}\quad 10\R{進数}{しん|すう}の\R{小数}{しょう|すう}の\R{位}{くらい}は、$\dfrac{1}{10}$・$\dfrac{1}{100}$…と、\R{右}{みぎ}へ1つ\R{進}{すす}むごとに $\dfrac{1}{10}$ になりました（2-1）。2\R{進数}{しん|すう}では、\R{右}{みぎ}へ1つ\R{進}{すす}むごとに $\dfrac{1}{2}$ になります。だから\R{小数}{しょう|すう}の\R{位}{くらい}は、$\dfrac{1}{2}$・$\dfrac{1}{4}$・$\dfrac{1}{8}$・$\dfrac{1}{16}$…です（$2^{-1}$、$2^{-2}$、…、2-7）。

\begin{bkcenter}
\begin{bktable}{colspec={Q[l,m,wd=34mm] Q[l,m,wd=56mm] Q[l,m,wd=26mm]}, rowsep=0.5pt, cells={font=\small}, row{1}={font=\sffamily\bfseries\small}}
10\R{進数}{しん|すう} & 2\R{進数}{しん|すう} & \\
$0.5=\dfrac{1}{2}$ & $0.1_{(2)}$ & \R{終}{お}わる \\
$0.75=\dfrac{1}{2}+\dfrac{1}{4}$ & $0.11_{(2)}$ & \R{終}{お}わる \\
$0.1$ & $0.000110011001100\cdots_{(2)}$ & \textbf{\R{終}{お}わらない} \\
\end{bktable}
\bktabcaption{0.1は、2\R{進数}{しん|すう}では「0011」がくり\R{返}{かえ}されて\R{終}{お}わらない。}
\end{bkcenter}

\textbf{0.1を2\R{進数}{しん|すう}にしてみる}\quad 2-17で\R{整数}{せい|すう}を2\R{進数}{しん|すう}にしたのと\R{同}{おな}じように、\R{大}{おお}きい\R{位}{くらい}から\R{順}{じゅん}に、\R{使}{つか}えたらひきます。

\begin{bkcenter}
\begin{bktable}{colspec={Q[c,m,wd=27mm] Q[l,m,wd=72mm] Q[c,m,wd=14mm]}, rowsep=1.6pt, colsep=3pt, cells={font=\small}, row{1}={font=\sffamily\bfseries\footnotesize}}
\R{位}{くらい} & \R{残}{のこ}りから\R{使}{つか}える？ & \R{数字}{すう|じ} \\
$\dfrac{1}{2}=0.5$ & \R{使}{つか}わない（\R{残}{のこ}り0.1より\R{大}{おお}きい） & 0 \\
$\dfrac{1}{4}=0.25$ & \R{使}{つか}わない & 0 \\
$\dfrac{1}{8}=0.125$ & \R{使}{つか}わない & 0 \\
$\dfrac{1}{16}=0.0625$ & \R{使}{つか}う：$0.1-0.0625=0.0375$ & 1 \\
$\dfrac{1}{32}=0.03125$ & \R{使}{つか}う：$0.0375-0.03125=0.00625$ & 1 \\
$\dfrac{1}{64}$、$\dfrac{1}{128}$ & \R{使}{つか}わない（\R{残}{のこ}り0.00625より\R{大}{おお}きい） & 0、0 \\
$\dfrac{1}{256}$、$\dfrac{1}{512}$ & \R{使}{つか}う、\R{使}{つか}う（\R{残}{のこ}りは0.000390625） & 1、1 \\
$\vdots$ & \R{残}{のこ}りはいつまでも0にならない & $\vdots$ \\
\end{bktable}
\bktabcaption{\R{残}{のこ}りが0にならないので、2\R{進数}{しん|すう}のけたが\R{終}{お}わらない。}
\end{bkcenter}

5\R{番目}{ばん|め}の\R{位}{くらい}のあとの\R{残}{のこ}り0.00625は、ちょうど $0.1\div16$ です。0.1を16でわった\R{数}{かず}が\R{残}{のこ}ったので、ここから\R{先}{さき}は、\R{位}{くらい}を4つずらして\R{同}{おな}じことがくり\R{返}{かえ}されます。だから「0011」がいつまでも\R{続}{つづ}くのです。

0.1は、10\R{進数}{しん|すう}の $\dfrac{1}{3}=0.333\cdots$ のように、2\R{進数}{しん|すう}ではどこまでも\R{続}{つづ}きます（2-9）。コンピュータはけたを\R{途中}{と|ちゅう}で\R{切}{き}って\R{覚}{おぼ}えるので、0.1はほんの\R{少}{すこ}しずれ、たし\R{算}{ざん}をくり\R{返}{かえ}すとずれが\R{見}{み}えてきます。

\bkhatenaref の\R{答}{こた}え：コンピュータは\textbf{まちがえたのではなく、0.1をぴったり\R{覚}{おぼ}えられない}のです。

\textbf{10\R{進数}{しん|すう}でたとえると}\quad 4けたまでしか\R{書}{か}けない\R{計算機}{けい|さん|き}では、$\dfrac{1}{3}$ は0.3333と\R{覚}{おぼ}えるしかありません。3\R{回}{かい}たすと $0.3333+0.3333+0.3333=0.9999$ で、1になりません。\R{本当}{ほん|とう}は $\dfrac{1}{3}\times3=1$ なのに、\R{覚}{おぼ}えたときに\R{切}{き}り\R{捨}{す}てた\R{分}{ぶん}が、たすたびにたまっていくのです。

\begin{motto}[2]
\textbf{2\R{倍}{ばい}していく\R{方法}{ほう|ほう}で、くり\R{返}{かえ}しを\R{見}{み}る}\par
10\R{進数}{しん|すう}の\R{小数}{しょう|すう}を10\R{倍}{ばい}すると、\R{小数点}{しょう|すう|てん}が\R{右}{みぎ}へ1つ\R{動}{うご}き、$\dfrac{1}{10}$ の\R{位}{くらい}の\R{数字}{すう|じ}が\R{整数}{せい|すう}の\R{位}{くらい}に\R{出}{で}てきました（2-1）。2\R{進数}{しん|すう}では、2\R{倍}{ばい}すると $\dfrac{1}{2}$ の\R{位}{くらい}の\R{数字}{すう|じ}が\R{整数}{せい|すう}の\R{位}{くらい}に\R{出}{で}てきます。だから、2\R{倍}{ばい}して\R{整数}{せい|すう}の\R{部分}{ぶ|ぶん}（0か1）を\R{書}{か}き、\R{小数}{しょう|すう}の\R{部分}{ぶ|ぶん}をまた2\R{倍}{ばい}する、をくり\R{返}{かえ}すと、2\R{進数}{しん|すう}のけたが\R{順}{じゅん}に\R{出}{で}てきます。
\begin{bkcenter}
\begin{bktable}{colspec={*{6}{Q[c,m,wd=17mm]}}, rowsep=0.5pt, colsep=2pt, cells={font=\small}, column{6}={bg=bkYellow!20}}
$0.1\times2$ & $0.2\times2$ & $0.4\times2$ & $0.8\times2$ & $0.6\times2$ & $0.2\times2$ \\
$=0.2$ & $=0.4$ & $=0.8$ & $=1.6$ & $=1.2$ & $=0.4$ \\
0 & 0 & 0 & 1 & 1 & 0 \\
\end{bktable}
\end{bkcenter}
6\R{回目}{かい|め}で、2\R{回目}{かい|め}と\R{同}{おな}じ「0.2×2」にもどりました。ここから\R{先}{さき}は、0、0、1、1、0、0、1、1、…と\R{同}{おな}じ\R{数字}{すう|じ}がくり\R{返}{かえ}されます。2-9で、わり\R{算}{ざん}のあまりが\R{前}{まえ}と\R{同}{おな}じになると\R{小数}{しょう|すう}がくり\R{返}{かえ}したのと、\R{同}{おな}じしくみです。
\end{motto}

\bksubsection{6}{\R{近}{ちか}い2つの\R{数}{かず}のひき\R{算}{ざん}と、ずれの\R{対策}{たい|さく}}

\leavevmode\label{bk:2-19:hikizan}コンピュータは、\R{小数}{しょう|すう}1つを64ビット（2-17）で\R{覚}{おぼ}えます。そのうち、\R{数字}{すう|じ}のならびを\R{覚}{おぼ}えるのに\R{使}{つか}うのは\R{約}{やく}53ビット\R{分}{ぶん}です。$2^{10}=1024$ はほぼ1000なので（2-7）、$2^{50}$ はほぼ $1000^5=10^{15}$ です。\R{実際}{じっ|さい}に $2^{53}=9007199254740992$ で16けたです。だから、コンピュータが\R{覚}{おぼ}えられる\R{数}{かず}は、10\R{進数}{しん|すう}にすると\textbf{\R{上}{うえ}から\R{約}{やく}16けた}までです。0.1のように2\R{進数}{しん|すう}で\R{終}{お}わらない\R{数}{かず}は、\R{上}{うえ}から17けためあたりに\R{小}{ちい}さなずれをもっています。

\R{大}{おお}きさのとても\R{近}{ちか}い2つの\R{数}{かず}をひき\R{算}{ざん}すると、どうなるでしょう。1000.1から1000をひいてみます。コンピュータの\R{中}{なか}の1000.1は、\R{上}{うえ}から17けためのあたりが\R{少}{すこ}しずれています。

\begin{bkfigure}
\begin{tikzpicture}[x=1mm, y=1mm]
  \node[anchor=east, font=\ttfamily\small, text=bkInk] at (78,12) {1000.1000000000000{\color{bkRed}\bfseries 227}};
  \node[anchor=east, font=\ttfamily\small, text=bkInk] at (78,6) {$-$\,1000.0000000000000000};
  \draw[bkInk, line width=0.6pt] (24,2.8) -- (79,2.8);
  \node[anchor=east, font=\ttfamily\small, text=bkInk] at (78,-1) {0.1000000000000{\color{bkRed}\bfseries 227}};
  \node[anchor=south east, font=\sffamily\footnotesize, text=bkInk] at (78,16) {\R{黒}{くろ}いけたは\R{正}{ただ}しい。{\color{bkRed}\R{赤}{あか}}いけたが\R{小}{ちい}さなずれ};
  \node[anchor=west, font=\sffamily\footnotesize, text=bkInk, align=left] at (82,12) {コンピュータの\R{中}{なか}の\\1000.1};
  \node[anchor=west, font=\sffamily\footnotesize, text=bkInk, align=left] at (82,-1) {\R{上}{うえ}の「1000」が\R{消}{き}えて、\\ずれ（\R{赤}{あか}）が\R{前}{まえ}に\R{出}{で}る};
\end{tikzpicture}
\bkfigcaption{1000.1－1000。\R{答}{こた}えの\R{正}{ただ}しいけたは、ひいて\R{消}{き}えた「1000」の\R{分}{ぶん}だけ\R{少}{すく}なくなる。}
\end{bkfigure}

ひき\R{算}{ざん}で\R{上}{うえ}の「1000」が\R{消}{き}えたので、\R{答}{こた}えの\R{正}{ただ}しいけたは4けたぶん\R{少}{すく}なくなり、\R{赤}{あか}いずれが\R{前}{まえ}に\R{出}{で}てきました。\R{大}{おお}きい\R{数}{かず}どうしほど、\R{消}{き}えるけたが\R{多}{おお}いので、ずれが\R{前}{まえ}に\R{出}{で}ます。

\begin{bkcenter}
\begin{bktable}{colspec={Q[l,m,wd=60mm] Q[l,m,wd=50mm]}, rowsep=0.5pt, cells={font=\small}, column{1-2}={font=\ttfamily\footnotesize}, row{1}={font=\sffamily\bfseries\small}}
プログラム & \R{実行}{じっ|こう}\R{結果}{けっ|か} \\
print(1.1 - 1) & 0.1000000000000000{\color{bkRed}\bfseries 9} \\
print(1000.1 - 1000) & 0.1000000000000{\color{bkRed}\bfseries 2274} \\
print(1000000000.1 - 1000000000) & 0.1000000{\color{bkRed}\bfseries 2384185791} \\
\end{bktable}
\bktabcaption{\R{本当}{ほん|とう}の\R{答}{こた}えはどれも0.1。\R{近}{ちか}い\R{大}{おお}きな\R{数}{かず}どうしほど、ずれ（\R{赤}{あか}）が\R{前}{まえ}に\R{出}{で}る。}
\end{bkcenter}

\textbf{コンピュータは\R{数}{かず}を\R{有限}{ゆう|げん}のけたしか\R{覚}{おぼ}えていないので、とても\R{近}{ちか}い2つの\R{数}{かず}をひき\R{算}{ざん}すると、\R{下}{した}のけたのずれが\R{目立}{め|だ}ちます}。「\R{有限}{ゆう|げん}」は「\R{限}{かぎ}りがある」という\R{意味}{い|み}です。\R{第}{だい}5\R{巻}{かん}の5-2では、このずれが\R{表計算}{ひょう|けい|さん}の\R{答}{こた}えに\R{出}{で}てきます。

\par\Needspace{8\baselineskip}\noindent\textbf{ずれの\R{対策}{たい|さく}}\par\nopagebreak
\begin{enumerate}[label={\bkcircled{\arabic*}}, leftmargin=2.2\zw, itemsep=1pt, topsep=2pt]
\item \textbf{\R{整数}{せい|すう}で\R{数}{かぞ}える}：\R{整数}{せい|すう}は2\R{進数}{しん|すう}でもぴったり\R{書}{か}けるので、ずれません（2-17）。\texttt{n = n + 1} を10\R{回}{かい}くり\R{返}{かえ}して10を\R{作}{つく}り、\R{最後}{さい|ご}に \texttt{n / 10} とすると、ぴったり \texttt{1.0} になります。
\item \textbf{ずれを\R{見}{み}こんで\R{比}{くら}べる}：\texttt{x == 1} ではなく、\texttt{0.999999 < x < 1.000001} のように「1にとても\R{近}{ちか}いか」を\R{比}{くら}べます（\R{不等号}{ふ|とう|ごう}、2-2）。
\item \textbf{\R{式}{しき}を\R{変形}{へん|けい}して、\R{近}{ちか}い\R{数}{かず}のひき\R{算}{ざん}をなくす}：\R{計算}{けい|さん}する\R{前}{まえ}に、\R{文字}{も|じ}の\R{式}{しき}のまま\R{整理}{せい|り}しておきます（5-2）。
\end{enumerate}

\begin{taisetsu}
プログラムの \texttt{=} は\R{代入}{だい|にゅう}：\R{右}{みぎ}を\R{計算}{けい|さん}して、\R{左}{ひだり}の\R{変数}{へん|すう}に\R{入}{い}れる。\R{変数}{へん|すう}の\R{中身}{なか|み}は\R{表}{ひょう}で\R{追}{お}う\\
0.1は、2\R{進数}{しん|すう}では\R{終}{お}わらないので、\R{覚}{おぼ}えると\R{少}{すこ}しずれる\\
とても\R{近}{ちか}い2つの\R{数}{かず}をひき\R{算}{ざん}すると、\R{下}{した}のけたのずれが\R{目立}{め|だ}つ
\end{taisetsu}

\begin{kotoba}
\small
\bkword{プログラム}{コンピュータへの\R{命令}{めい|れい}を、\R{順番}{じゅん|ばん}にならべたもの。}
\bkword{\R{変数}{へん|すう}・\R{代入}{だい|にゅう}}{\R{数}{かず}を\R{入}{い}れておく\R{名前}{な|まえ}つきの\R{箱}{はこ}が\R{変数}{へん|すう}。\R{箱}{はこ}に\R{数}{かず}を\R{入}{い}れることが\R{代入}{だい|にゅう}。}
\bkword{\R{条件分岐}{じょう|けん|ぶん|き}}{\R{条件}{じょう|けん}が\R{正}{ただ}しいかどうかで、\R{実行}{じっ|こう}する\R{命令}{めい|れい}を\R{変}{か}えること。}
\end{kotoba}

\bkkigocards{\bkkigocard{\ttfamily ==}{イコール・イコール}{プログラムで、\R{左}{ひだり}と\R{右}{みぎ}が\R{等}{ひと}しいかを\R{比}{くら}べる\R{記号}{き|ごう}。\texttt{=} は\R{代入}{だい|にゅう}。}{equal to}{2-19}}
{\bkkigocard{\ttfamily >=}{だいなりイコール}{プログラムで ≧ のかわりに\R{使}{つか}う\R{記号}{き|ごう}。\R{左}{ひだり}が\R{右}{みぎ}\R{以上}{い|じょう}。}{greater than or equal to}{2-19}}

%% ---------------- 例題 ----------------
\bksection{reidai}{\R{例題}{れい|だい}でたしかめよう}

\begin{reidai}[\R{代入}{だい|にゅう}の\R{順番}{じゅん|ばん}]
\noindent\begin{minipage}[t]{0.62\linewidth}
\R{右}{みぎ}のプログラムで\R{表示}{ひょう|じ}される2つの\R{数}{かず}を\R{答}{こた}えましょう。（\texttt{print(a, b)} は、a と b をならべて\R{表示}{ひょう|じ}します。）\par
\kangaekata \R{上}{うえ}から1\R{行}{ぎょう}ずつ、\R{箱}{はこ}の\R{中身}{なか|み}を\R{表}{ひょう}に\R{書}{か}く。b は、b を\R{計算}{けい|さん}した\R{時}{とき}の a を\R{使}{つか}うので、あとで a が\R{変}{か}わっても b は\R{変}{か}わらない。
\end{minipage}\hfill
\begin{minipage}[t]{0.34\linewidth}
\vspace{2pt}
\begin{bkiitcode}
a = 4
b = a * 3
a = a + 1
print(a, b)
\end{bkiitcode}
\end{minipage}\par
\shiki\begin{bkiittrace}[t]{colspec={Q[l,m,wd=30mm] Q[c,m,wd=10mm] Q[c,m,wd=10mm] Q[l,m,wd=46mm]}}
\R{実行}{じっ|こう}した\R{行}{ぎょう} & a & b & \R{計算}{けい|さん} \\
\texttt{a = 4} & 4 & － & \\
\texttt{b = a * 3} & 4 & 12 & $4\times3=12$ \\
\texttt{a = a + 1} & 5 & 12 & $4+1=5$。b は12のまま \\
\texttt{print(a, b)} & 5 & 12 & 「5\enspace 12」と\R{表示}{ひょう|じ} \\
\end{bkiittrace}\par
\kotae 5\quad 12
\end{reidai}

\begin{jibun}
\texttt{x = 10}、\texttt{y = x - 3}、\texttt{x = y * 2}、\texttt{print(x, y)} の\R{順}{じゅん}に\R{実行}{じっ|こう}すると、\R{何}{なに}が\R{表示}{ひょう|じ}されますか。\R{変数}{へん|すう}の\R{表}{ひょう}を\R{書}{か}いて\R{答}{こた}えましょう。
\bkjans{y は $10-3=7$、x は $7\times2=14$。\R{表示}{ひょう|じ}は「14\nobreak\quad\nobreak 7」}
\end{jibun}

\begin{reidai}[くり\R{返}{かえ}し]
\noindent\begin{minipage}[t]{0.62\linewidth}
\R{右}{みぎ}のプログラムで\R{表示}{ひょう|じ}される\R{数}{かず}は？\par
\kangaekata \texttt{range(5)} で、\R{下}{さ}げた\R{行}{ぎょう}を5\R{回}{かい}。\texttt{print} は\R{最後}{さい|ご}に1\R{回}{かい}。\par
\shiki s は 0→3→6→9→12→15\par
\kotae 15
\end{minipage}\hfill
\begin{minipage}[t]{0.34\linewidth}
\vspace{2pt}
\begin{bkiitcode}
s = 0
for i in range(5):
    s = s + 3
print(s)
\end{bkiitcode}
\end{minipage}
\end{reidai}


\begin{reidai}[もし～なら]
\noindent\begin{minipage}[t]{0.62\linewidth}
\R{右}{みぎ}のプログラムで\R{表示}{ひょう|じ}されることばは？\par
\kangaekata \texttt{t >= 25} は「t≧25」。\R{正}{ただ}しいか\R{調}{しら}べる。\par
\shiki $28\geqq25$ は\R{正}{ただ}しいので、\texttt{if} のすぐ\R{下}{した}を\R{実行}{じっ|こう}\par
\kotae なつび
\end{minipage}\hfill
\begin{minipage}[t]{0.34\linewidth}
\vspace{2pt}
\begin{bkiitcode}
t = 28
if t >= 25:
    print("なつび")
else:
    print("ふつう")
\end{bkiitcode}
\end{minipage}
\end{reidai}

\begingroup\tcbset{equal height group=bkiitjA}%
\noindent\begin{minipage}[t]{0.49\linewidth}
\begin{jibun}
\texttt{n = 1} のあと \texttt{n = n * 2} を4\R{回}{かい}くり\R{返}{かえ}し、\texttt{print(n)}。\R{何}{なに}が\R{表示}{ひょう|じ}される？
\bkjans{n は 1→2→4→8→16。16}
\end{jibun}
\end{minipage}\hfill
\begin{minipage}[t]{0.49\linewidth}
\begin{jibun}
\texttt{score = 55}。もし \texttt{score >= 60} なら「ごうかく」、でなければ「もういちど」と\R{表示}{ひょう|じ}。\R{何}{なに}が\R{出}{で}る？
\bkjans{$55\geqq60$ は\R{正}{ただ}しくないので、もういちど}
\end{jibun}
\end{minipage}\par
\endgroup


\begin{reidai}[くり\R{返}{かえ}しのあとで、もし～なら]
\noindent\begin{minipage}[t]{0.55\linewidth}
\R{右}{みぎ}のプログラムで\R{表示}{ひょう|じ}されることばを\R{答}{こた}えましょう。\par
\kangaekata w を\R{求}{もと}めてから、\texttt{if} を\R{調}{しら}べる。\par
\shiki w は $20-3\times4=8$。$8\geqq10$ は\R{正}{ただ}しくないので、\texttt{else} の\R{下}{した}を\R{実行}{じっ|こう}\par
\kotae たりない
\end{minipage}\hfill
\begin{minipage}[t]{0.41\linewidth}
\vspace{2pt}
\begin{bkiitcode}
w = 20
for i in range(4):
    w = w - 3
if w >= 10:
    print("じゅうぶん")
else:
    print("たりない")
\end{bkiitcode}
\end{minipage}
\end{reidai}


\begin{reidai}[2\R{進数}{しん|すう}で\R{終}{お}わる\R{小数}{しょう|すう}]
0.375は、2\R{進数}{しん|すう}で\R{終}{お}わる\R{小数}{しょう|すう}ですか。\R{終}{お}わるなら2\R{進数}{しん|すう}で\R{書}{か}きましょう。\par
\kangaekata \R{大}{おお}きい\R{位}{くらい}（$\dfrac{1}{2}=0.5$、$\dfrac{1}{4}=0.25$、$\dfrac{1}{8}=0.125$、…）から\R{順}{じゅん}に、\R{使}{つか}えたらひく。\R{残}{のこ}りが0になれば\R{終}{お}わる。\par
\shiki 0.5は\R{使}{つか}わない（0.375より\R{大}{おお}きい）\par
\hspace*{2.6\zw}0.25を\R{使}{つか}う：$0.375-0.25=0.125$\par
\hspace*{2.6\zw}0.125を\R{使}{つか}う：$0.125-0.125=0$\quad \R{残}{のこ}りが0になった\par
\kotae \R{終}{お}わる。$\dfrac{1}{2}$ の\R{位}{くらい}が0、$\dfrac{1}{4}$ と $\dfrac{1}{8}$ の\R{位}{くらい}が1で、$0.011_{(2)}$
\end{reidai}

\begingroup\tcbset{equal height group=bkiitjB}%
\noindent\begin{minipage}[t]{0.49\linewidth}
\begin{jibun}
\texttt{p = 50} のあと \texttt{p = p + 20} を3\R{回}{かい}。そのあと、もし \texttt{p >= 100} なら「OK」、でなければ「NG」と\R{表示}{ひょう|じ}。\R{何}{なに}が\R{出}{で}る？
\bkjans{p は $50+20\times3=110$。$110\geqq100$ なので OK}
\end{jibun}
\end{minipage}\hfill
\begin{minipage}[t]{0.49\linewidth}
\begin{jibun}
0.625は、2\R{進数}{しん|すう}で\R{終}{お}わる\R{小数}{しょう|すう}ですか。\R{終}{お}わるなら2\R{進数}{しん|すう}で\R{書}{か}きましょう。
\bkjans{$0.625-0.5=0.125$、$\dfrac{1}{4}$ は\R{使}{つか}わない、$0.125-0.125=0$。\R{終}{お}わる。$0.101_{(2)}$}
\end{jibun}
\end{minipage}\par
\endgroup


\begin{reidai}[\R{単位}{たん|い}の\R{換算}{かん|さん}をプログラムで]
\noindent\begin{minipage}[t]{0.6\linewidth}
\R{右}{みぎ}のプログラムは、1\R{日}{にち}が\R{何秒}{なん|びょう}かを\R{計算}{けい|さん}します。それぞれの\R{行}{ぎょう}の24や60は\R{何}{なに}を\R{表}{あらわ}していますか。また、\R{表示}{ひょう|じ}される\R{数}{かず}を\R{答}{こた}えましょう。\par
\kangaekata 1\R{日}{にち}は24\R{時間}{じ|かん}（\qty{24}{h/d}）、1\R{時間}{じ|かん}は60\R{分}{ぷん}（\qty{60}{min/h}）、1\R{分}{ぷん}は60\R{秒}{びょう}（\qty{60}{s/min}）。dはday（\R{日}{にち}）の\R{頭文字}{かしら|も|じ}です（2-10）。
\end{minipage}\hfill
\begin{minipage}[t]{0.36\linewidth}
\vspace{2pt}
\begin{bkiitcode}
d = 1
h = d * 24
m = h * 60
s = m * 60
print(s)
\end{bkiitcode}
\end{minipage}\par
\shiki \R{単位}{たん|い}つきの\R{式}{しき}で\R{書}{か}くと、1\R{行}{ぎょう}ずつ\R{単位}{たん|い}が\R{約分}{やく|ぶん}されていきます。
\begin{align*}
\qty{1}{d}\times\qty{24}{h/d}&=\qty{24}{h}\\
\qty{24}{h}\times\qty{60}{min/h}&=\qty{1440}{min}\\
\qty{1440}{min}\times\qty{60}{s/min}&=\qty{86400}{s}
\end{align*}
\hspace*{2.6\zw}h の\R{行}{ぎょう}の24は\qty{24}{h/d}、m の\R{行}{ぎょう}の60は\qty{60}{min/h}、s の\R{行}{ぎょう}の60は\qty{60}{s/min}\par
\kotae 86400（1\R{日}{にち}は\qty{86400}{s}）
\end{reidai}

\begin{jibun}
1\R{時間}{じ|かん}が\R{何}{なん}ミリ\R{秒}{びょう}（ms）かを\R{計算}{けい|さん}するプログラムを、\texttt{h = 1} から\R{始}{はじ}めて3\R{行}{ぎょう}で\R{書}{か}き、\R{表示}{ひょう|じ}される\R{数}{かず}を\R{答}{こた}えましょう。1\R{秒}{びょう}は1000ミリ\R{秒}{びょう}（\qty{1000}{ms/s}、2-10）です。
\bkjans{\texttt{h = 1}、\texttt{s = h * 3600}、\texttt{ms = s * 1000}、\texttt{print(ms)}。$\qty{1}{h}\times\qty{3600}{s/h}\times\qty{1000}{ms/s}=\qty{3600000}{ms}$ で3600000}
\end{jibun}

\clearpage
%% ---------------- まちがえやすいところ・物理工学 ----------------
\bksection{tsumazuki}{まちがえやすいところをチェック}

\begin{tsumazuki}{\texttt{x = x + 1} を\R{方程式}{ほう|てい|しき}だと\R{思}{おも}う}
\bkng{\texttt{x = x + 1} は、x と x＋1 が\R{等}{ひと}しいという\R{意味}{い|み}だから、おかしい}
\bkok{\texttt{x = x + 1} は「\R{今}{いま}の x に1をたして、x に\R{入}{い}れなおす」という\R{命令}{めい|れい}}
プログラムの \texttt{=} は\R{代入}{だい|にゅう}です。\R{等}{ひと}しいかを\R{比}{くら}べるときは \texttt{==} を\R{使}{つか}います。
\end{tsumazuki}

\begin{tsumazuki}{\R{下}{さ}げた\R{行}{ぎょう}と\R{下}{さ}げない\R{行}{ぎょう}を\R{区別}{く|べつ}しない}
\bkng{\texttt{for i in range(3):} の\R{下}{した}に\R{下}{さ}げずに \texttt{print(x)} → 3\R{回}{かい}\R{表示}{ひょう|じ}される}
\bkok{\R{下}{さ}げずに\R{書}{か}いた \texttt{print(x)} は、くり\R{返}{かえ}しが\R{終}{お}わってから1\R{回}{かい}だけ\R{実行}{じっ|こう}される}
Pythonでは、\R{右}{みぎ}へ\R{下}{さ}げた\R{行}{ぎょう}だけがくり\R{返}{かえ}しの\R{中身}{なか|み}です。
\end{tsumazuki}

\begin{tsumazuki}{\R{小数}{しょう|すう}を \texttt{==} で\R{比}{くら}べる}
\bkng{0.1を10\R{回}{かい}たした x を \texttt{if x == 1:} で\R{比}{くら}べる（「1ではない」と\R{判断}{はん|だん}される）}
\bkok{\texttt{if 0.999999 < x < 1.000001:} のように、ずれを\R{見}{み}こんで\R{比}{くら}べる}
x は0.9999999999999999なので、\texttt{x == 1} は\R{正}{ただ}しくないことになってしまいます。
\end{tsumazuki}
\begin{butsuri}
\textbf{\R{第}{だい}5\R{巻}{かん}の「\R{瞬間}{しゅん|かん}の\R{速}{はや}さ」}\par
5-2では、\R{落}{お}ちるボールの1\R{秒後}{びょう|ご}の\R{瞬間}{しゅん|かん}の\R{速}{はや}さ（\R{本当}{ほん|とう}は\qty{10}{m/s}）を、とても\R{短}{みじか}い\R{時間}{じ|かん} $h$ \R{秒}{びょう}の\R{間}{あいだ}の\R{平均}{へい|きん}の\R{速}{はや}さ $\dfrac{5(1+h)^2-5}{h}$ で\R{求}{もと}めます。$h$ を\R{小}{ちい}さくするほど\R{本当}{ほん|とう}の\R{速}{はや}さに\R{近}{ちか}づく…はずです。

\noindent\begin{minipage}[c]{0.48\linewidth}
\begin{bkiitcode}
h = 0.00000001
x = 5 * (1 + h) ** 2
print((x - 5) / h)
\end{bkiitcode}
{\footnotesize\sffamily \texttt{**} は\R{累乗}{るい|じょう}（2-7）。\texttt{(1 + h) ** 2} は $(1+h)^2$\par}
\end{minipage}\hfill
\begin{minipage}[c]{0.48\linewidth}\centering
\begin{bktable}{colspec={Q[c,m,wd=12mm] Q[l,m,wd=38mm]}, rowsep=0.3pt, colsep=4pt, cells={font=\footnotesize}, row{1}={font=\sffamily\bfseries\footnotesize}}
$h$ & \R{実行}{じっ|こう}\R{結果}{けっ|か} \\
0.1 & \texttt{10.500000000000007} \\
0.001 & \texttt{10.00499999999871} \\
$10^{-5}$ & \texttt{10.000050000069649} \\
$10^{-8}$ & \texttt{9.99999993922529} \\
$10^{-10}$ & \texttt{{\color{bkRed}\bfseries 10.00000082740371}} \\
$10^{-12}$ & \texttt{{\color{bkRed}\bfseries 10.00088900582341}} \\
\end{bktable}
\end{minipage}

\par\addvspace{4pt}
$h$ を $10^{-8}$ より\R{小}{ちい}さくすると、\R{逆}{ぎゃく}に10からはなれていきます。$10^{-8}$ でも、\R{本当}{ほん|とう}の\R{値}{あたい}10.00000005より\R{小}{ちい}さくなっています。x と5がとても\R{近}{ちか}く、ひき\R{算}{ざん} \texttt{x - 5} で\R{下}{した}のけたのずれが\R{目立}{め|だ}つのです。そのうえ、とても\R{小}{ちい}さい $h$ でわるので、ずれが\R{大}{おお}きくなります。5-2では\R{式}{しき}を $10+5h$ に\R{変形}{へん|けい}して、ひき\R{算}{ざん}そのものをなくします。

\par\addvspace{3pt}\textbf{\R{時計}{とけい}やお\R{金}{かね}は\R{整数}{せい|すう}で}\par\noindent コンピュータの\R{時計}{とけい}は「\R{何回目}{なん|かい|め}の1ミリ\R{秒}{びょう}か」を、お\R{金}{かね}は\R{円}{えん}を、\R{整数}{せい|すう}で\R{数}{かぞ}えます。1\R{日}{にち}は $\qty{86400}{s}\times\qty{1000}{ms/s}=\qty{86400000}{ms}$ なので、1\R{日}{にち}の\R{時刻}{じ|こく}は0から86399999までの\R{整数}{せい|すう}で\R{表}{あらわ}せます。\R{大学}{だい|がく}の「\R{数値計算}{すう|ち|けい|さん}」では、ずれの\R{小}{ちい}さい\R{計算}{けい|さん}の\R{手順}{て|じゅん}を\R{考}{かんが}えます。
\end{butsuri}

\begin{butsuri}
\textbf{\R{冷蔵庫}{れい|ぞう|こ}の\R{温度}{おん|ど}を\R{守}{まも}るプログラム}\par
\R{冷蔵庫}{れい|ぞう|こ}の\R{中}{なか}の\R{小}{ちい}さなコンピュータは、くり\R{返}{かえ}しと「もし～なら」で\R{温度}{おん|ど}を\qty{5}{\degreeCelsius}くらいに\R{保}{たも}っています。1\R{分}{ぷん}ごとに\R{温度}{おん|ど}センサー（2-17）で\R{温度}{おん|ど}をはかり、\qty{5}{\degreeCelsius}より\R{高}{たか}ければ\R{冷}{ひ}やす\R{機械}{き|かい}を\R{動}{うご}かし、そうでなければ\R{止}{と}めます。

\par\addvspace{6pt}\noindent\begin{minipage}[c]{0.5\linewidth}
\begin{bkiitcode}
for i in range(1440):
    t = sensor()
    if t > 5:
        cool\_on()
    else:
        cool\_off()
    wait\_1min()
\end{bkiitcode}
{\footnotesize\sffamily \texttt{sensor()} は「\R{温度}{おん|ど}をはかる」、\texttt{cool\_on()} は「\R{冷}{ひ}やす」、\texttt{wait\_1min()} は「1\R{分}{ぷん}\R{待}{ま}つ」という\R{命令}{めい|れい}（\R{作}{つく}った\R{会社}{かい|しゃ}が\R{用意}{よう|い}したもの）。\par}
\end{minipage}\hfill
\begin{minipage}[c]{0.46\linewidth}\centering
\begin{tikzpicture}
\begin{axis}[bkplot, width=60mm, height=42mm, xmin=0, xmax=40, ymin=3, ymax=6.4, xtick={0,10,20,30,40}, ytick={4,5,6},
  xlabel={\R{分}{ふん}}, ylabel={℃}]
  \addplot[bkthin, dashed, domain=0:40] {5};
  \addplot[bkcurve, line width=1pt] coordinates {(0,4.2) (6,5.1) (9,4.3) (15,5.1) (18,4.3) (24,5.1) (27,4.3) (33,5.1) (36,4.3) (40,4.9)};
\end{axis}
\end{tikzpicture}
\par{\footnotesize\sffamily \R{止}{と}めると\R{上}{あ}がり、\R{冷}{ひ}やすと\R{下}{さ}がる。\par}
\end{minipage}

\par\addvspace{4pt}
1440は、1\R{日}{にち}が\R{何分}{なん|ぷん}かを\R{表}{あらわ}す\R{数}{かず}です（1\R{分}{ぷん}に1\R{回}{かい}ずつ、1\R{日分}{にち|ぶん}くり\R{返}{かえ}す）。
\[ \qty{1}{d}\times\qty{24}{h/d}\times\qty{60}{min/h}=\qty{1440}{min} \]
このように、はかった\R{結果}{けっ|か}を\R{見}{み}て\R{機械}{き|かい}の\R{動}{うご}きを\R{決}{き}めることを\textbf{\R{制御}{せい|ぎょ}}といいます。エアコンや\R{自動車}{じ|どう|しゃ}のエンジン、ロボットも、\R{同}{おな}じ\R{考}{かんが}え\R{方}{かた}のプログラムで\R{動}{うご}いています。\R{大学}{だい|がく}の「\R{制御工学}{せい|ぎょ|こう|がく}」では、オン・オフだけでなく、\R{冷}{ひ}やす\R{強}{つよ}さを\R{少}{すこ}しずつ\R{変}{か}える\R{上手}{じょう|ず}な\R{制御}{せい|ぎょ}を\R{学}{まな}びます。
\end{butsuri}


\clearpage
%% ---------------- 章末問題 ----------------
\bkendproblems
\bkproblemgroup{\R{基本}{き|ほん}}{2}

\bkproblem \R{次}{つぎ}のプログラムを\R{実行}{じっ|こう}すると、\R{何}{なに}が\R{表示}{ひょう|じ}されますか。
\par\nobreak\vspace{4pt}
\noindent\begin{minipage}[t]{0.31\linewidth}
（1）
\begin{bkiitcode}
a = 7
b = a * 2
a = a + b
print(a)
\end{bkiitcode}
\end{minipage}\hfill
\begin{minipage}[t]{0.33\linewidth}
（2）
\begin{bkiitcode}
s = 100
for i in range(4):
    s = s - 15
print(s)
\end{bkiitcode}
\end{minipage}\hfill
\begin{minipage}[t]{0.31\linewidth}
（3）
\begin{bkiitcode}
n = 3
for i in range(3):
    n = n * 2
print(n)
\end{bkiitcode}
\end{minipage}
\bkpans{（1）\nobreak b は14、a は $7+14=21$。21\quad（2）\nobreak $100-15\times4=40$\quad（3）\nobreak $3\times2\times2\times2=24$}

\bkproblem \R{次}{つぎ}のプログラムを\R{実行}{じっ|こう}すると、\R{何}{なに}が\R{表示}{ひょう|じ}されますか。
\par\nobreak\vspace{4pt}
\noindent\begin{minipage}[t]{0.4\linewidth}
\begin{bkiitcode}
x = 12
if x >= 15:
    print("おおきい")
else:
    print("ちいさい")
\end{bkiitcode}
\end{minipage}
\bkpans{$12\geqq15$ は\R{正}{ただ}しくないので、ちいさい}

\bkproblem \R{次}{つぎ}の\R{小数}{しょう|すう}のうち、2\R{進数}{しん|すう}で\R{終}{お}わるものをすべて\R{選}{えら}び、2\R{進数}{しん|すう}で\R{書}{か}きましょう。\quad ア\enspace 0.5\quad イ\enspace 0.125\quad ウ\enspace 0.3\quad エ\enspace 0.75
\bkpans{ア $0.1_{(2)}$、イ $0.001_{(2)}$、エ $0.11_{(2)}$}

\bkproblem プログラムの \texttt{x = x + 1} と、\R{方程式}{ほう|てい|しき} $x=x+1$ のちがいを\R{説明}{せつ|めい}しましょう。
\bkpans{プログラムは「x に1をたして x に\R{入}{い}れなおす」\R{代入}{だい|にゅう}。\R{方程式}{ほう|てい|しき}は x と x＋1 が\R{等}{ひと}しいという\R{意味}{い|み}で、\R{解}{かい}はない}

\bkproblem \R{次}{つぎ}のプログラムで\R{表示}{ひょう|じ}される\R{数}{かず}を\R{答}{こた}えましょう。また、それは\R{何}{なに}を\R{計算}{けい|さん}していますか。1\R{週間}{しゅう|かん}は7\R{日}{にち}です。
\par\nobreak\vspace{4pt}
\noindent\begin{minipage}[t]{0.4\linewidth}
\begin{bkiitcode}
w = 3
d = w * 7
h = d * 24
print(h)
\end{bkiitcode}
\end{minipage}
\bkpans{3\R{週間}{しゅう|かん}は\qty{21}{d}（7\R{日}{にち}×3）、\qty{21}{d}×\qty{24}{h/d}＝\qty{504}{h}。504。3\nobreak\R{週間}{しゅう|かん}が\R{何時間}{なん|じ|かん}か（\R{日}{にち}にしてから\R{時間}{じ|かん}にしている）}

\clearpage
\bkproblemgroup{チャレンジ}{2}

\bkproblem \R{次}{つぎ}のプログラムを\R{実行}{じっ|こう}すると、\R{何}{なに}が\R{表示}{ひょう|じ}されますか。
\par\nobreak\vspace{4pt}
\noindent\begin{minipage}[t]{0.45\linewidth}
\begin{bkiitcode}
m = 1000
for i in range(10):
    m = m + 500
if m >= 5000:
    print("たまった")
else:
    print("まだ")
\end{bkiitcode}
\end{minipage}
\bkpans{m は $1000+500\times10=6000$。$6000\geqq5000$ なので、たまった}

\bkproblem 0.25を4\R{回}{かい}たした x について、\texttt{if x == 1:} で\R{比}{くら}べると、「\R{等}{ひと}しい」と\R{判断}{はん|だん}されます。0.1を10\R{回}{かい}たしたときとちがうのはなぜですか。
\bkpans{0.25は $\dfrac{1}{4}$ で、2\R{進数}{しん|すう}で $0.01_{(2)}$ と\R{終}{お}わるので、ぴったり\R{覚}{おぼ}えられ、ずれが\R{出}{で}ないから}

\bkproblem 0.1を10\R{回}{かい}たして1を\R{作}{つく}りたいとき、ずれが\R{出}{で}ないようにするには、どう\R{計算}{けい|さん}すればよいですか。1つ\R{説明}{せつ|めい}しましょう。
\bkpans{（\R{例}{れい}）1を10\R{回}{かい}たして\R{整数}{せい|すう}の10を\R{作}{つく}り、\R{最後}{さい|ご}に10でわる}

\begin{furikaeri}
\begin{checklist}
\item \R{変数}{へん|すう}と\R{代入}{だい|にゅう}のプログラムの\R{結果}{けっ|か}を\R{予想}{よ|そう}できる。\hfill{\footnotesize\sffamily\color{bkTheme}→\pageref{bk:2-19:dainyu}ページ}
\item くり\R{返}{かえ}しと「もし～なら」のプログラムを\R{読}{よ}める。\hfill{\footnotesize\sffamily\color{bkTheme}→\pageref{bk:2-19:kurikaesi}ページ}
\item \R{小数}{しょう|すう}のずれのわけと\R{対策}{たい|さく}を\R{説明}{せつ|めい}できる。\hfill{\footnotesize\sffamily\color{bkTheme}→\pageref{bk:2-19:zure}ページ}
\end{checklist}
\end{furikaeri}

\bkfukidashi{delta}{コンピュータは\R{計算}{けい|さん}が\R{速}{はや}いけれど、\R{数}{かず}を\R{覚}{おぼ}えるけたには\R{限}{かぎ}りがあるんだ。\R{式}{しき}で\R{考}{かんが}える\R{力}{ちから}と\R{組}{く}み\R{合}{あ}わせて\R{使}{つか}おう。}
